\section{Introduction}
The standard cosmological model describes a wide range of observations while leaving open questions concerning initial conditions, singularities, vacuum energy, and the global structure of cosmological history \cite{planck2018,desi2025dr2}. Cyclic and bouncing cosmologies ask whether the observed expanding phase can be embedded in a nonsingular recurrent history \cite{ashtekar2011,lehners2008}. Such a construction is more restrictive than drawing a periodic scale factor: it must follow from declared equations, satisfy the gravitational constraints, remain regular and stable, and yield observables.

The purpose of \SCPC{} is not to infer cyclicity from suggestive numerical patterns. It is to define a model hierarchy in which each proposed mechanism is covariant, dimensionally explicit, numerically verifiable, perturbatively stable, and confrontable with public likelihoods. This paper performs three bounded tasks. First, it derives and tests the conservative four-dimensional canonical baseline, including exact local turning-point conditions and an explicit numerical non-cycle result for the configured state. Second, it validates the background-observable and likelihood machinery by reproducing the public DESI DR2 flat-$\Lambda$CDM BAO and BBN posterior constraints from the official chain release. Third, it formalizes the proposed ``simultaneous layers'' geometry as an $S^3$ foliation and implements the maximum higher-dimensional calculation justified by the currently declared physics: a symmetry-reduced scalar test field on fixed $\mathbb R\times S^4$.

The central negative result is scientifically useful. A canonical scalar obeys the null energy condition, its configured positive potential offset produces monotonic accelerated expansion, and no cyclic solution has yet been established. Neither an embedding coordinate nor a numerical return path is identified with time. A self-gravitating layered model remains undefined until a higher-dimensional action, boundary conditions, constraint formulation, and matter couplings are supplied. Likewise, the DESI calculation in this paper is a standard-cosmology reproduction benchmark for the analysis pipeline, not an observational fit of \SCPC{} itself.
